\clearpage
\phantomsection
\addcontentsline{toc}{chapter}{Executive Summary}
\chapter*{Executive Summary}

Customer churn, defined as the discontinuation of service by existing users, is one of the most important business risks for FinTech platforms. It lowers recurring revenue, increases customer acquisition cost, and weakens long-term growth. Prior studies note that retaining an existing customer is significantly more cost-effective than acquiring a new one, which makes early churn identification a practical and strategic priority \cite{systematicreview,bhattacharjee2026}. This project therefore focuses on developing a machine learning based churn prediction framework for a FinTech platform.

The proposed work assumes an anonymized customer dataset of approximately 10{,}000 records with demographic, transactional, account-level, and support-related variables. Typical features include age, gender, region, balance, credit score, tenure, transaction count, complaints, subscription type, and account activity. The project pipeline begins with data preprocessing, including removal of irrelevant identifiers, missing-value handling, outlier treatment, categorical encoding, and scaling of numerical variables. Because churn is usually a minority class, imbalance handling methods such as SMOTE and class-weight adjustment are also incorporated.

Feature engineering is used to improve predictive performance by creating behavior-focused variables such as average transactions per month, complaint rate, and tenure buckets. Multiple classification models are then trained and compared, namely Logistic Regression, Decision Tree, Random Forest, and XGBoost. These algorithms are selected to provide a meaningful balance between interpretability and predictive strength. Hyperparameters are tuned through grid search and cross-validation, and performance is evaluated using accuracy, precision, recall, F1-score, ROC-AUC, and confusion matrix analysis.

The example comparative results in this report follow the trend reported in literature: ensemble methods outperform simpler baselines, with XGBoost producing the strongest overall performance in terms of churn-class recall and F1-score \cite{imani2024,bhattacharjee2026}. The final system is designed not only to classify customers as churners or non-churners, but also to assign operational risk categories such as High, Medium, and Low. These predictions can be integrated into a web dashboard for business teams, allowing managers to identify at-risk customers and prioritize interventions efficiently.

This report has been structured as a complete academic major project document. It includes title and certificate pages, acknowledgements, executive summary, abstract, literature review, system analysis, design, implementation, testing, results, conclusions, references, and appendices. Architecture, UML, DFD, and project timeline diagrams are included in editable Mermaid code form, and placeholder figures have been inserted for experimental plots that can be replaced later by actual outputs. Where exact dataset or institutional details were unspecified, explicit assumptions have been stated. The document is therefore ready for compilation, review, and final formatting for submission.
